% ---- prefix ----
\documentclass[11pt,letterpaper]{article}
\usepackage[margin=1in]{geometry}
\usepackage{amsmath,amssymb,amsthm,mathtools,mathrsfs}
\usepackage[T1]{fontenc}
\usepackage{lmodern,microtype}
\usepackage[dvipsnames]{xcolor}
\usepackage{booktabs,array}
\usepackage{enumitem}
\usepackage[colorlinks=true,linkcolor=MidnightBlue,citecolor=MidnightBlue,urlcolor=MidnightBlue]{hyperref}
\numberwithin{equation}{section}
\newtheorem{theorem}{Theorem}[section]
\newtheorem{lemma}[theorem]{Lemma}
\newtheorem{proposition}[theorem]{Proposition}
\newtheorem{corollary}[theorem]{Corollary}
\theoremstyle{definition}
\newtheorem{definition}[theorem]{Definition}
\newtheorem{conjecture}[theorem]{Conjecture}
\theoremstyle{remark}
\newtheorem{remark}[theorem]{Remark}
\DeclareMathOperator{\sgn}{sgn}
\DeclareMathOperator{\rot}{rot}
\DeclareMathOperator{\lk}{lk}
\DeclareMathOperator{\Tr}{Tr}
\DeclareMathOperator{\diag}{diag}
\DeclareMathOperator{\slk}{sl}
\newcommand{\R}{\mathbb R}
\newcommand{\Z}{\mathbb Z}
\newcommand{\Dcal}{\mathcal D}
\newcommand{\vseven}[1]{\texttt{[eq::#1]}}
\newcommand{\coeff}[1]{\bigl[#1\bigr]}
\setlength{\emergencystretch}{2em}
\setlist{topsep=4pt,itemsep=3pt}
\title{Marked polygon currents and the HOMFLY corner\\
\large A direct sewing argument for the GLSSW state sum}
\author{}
\date{13 September 2026}
\begin{document}
\maketitle
\begin{abstract}
We reorganize the finite single-minus recursion of GLSSW into effective
polygon vertices. Separating-smoothing inversion defines a primitive
\(\Gamma_A\). An interval-current calculation makes this primitive vanish
on every nonuniform polygon. On a uniform polygon, an open current with a
marked, auxiliary closing edge has precisely the descending skein recursion
of a framed HOMFLY coefficient. This identifies its primitive with the
corner \([a^{1-m-|r|}z^0]H^+\) and reconstructs the uniform-decomposition
state sum. The argument uses GLSSW root independence, finite cut-list
identities, and elementary knot skein theory; it does not use a
wall-crossing uniqueness theorem, a transverse degree bound, or the
state-sum equality as an input. We then give the Wilson-loop normalization
and distinguish this algebraic sewing from a conjectural realization in
topological string field theory.
\end{abstract}
{\footnotesize\tableofcontents}
\clearpage

\section{Statement, conventions, and inputs}
\label{sec:statement}

Let \(P=(q_0,\ldots,q_{n-1})\) be an oriented polygon immersed in an
oriented affine plane. Its external momentum edges are
\(\widetilde\lambda_i=q_{i+1}-q_i\), with cyclic indices. A left turn means
\(\det(\widetilde\lambda_{i-1},\widetilde\lambda_i)>0\).
Write \(\ell(P)\) for the number of left turns, and \(r(P)\) for the signed
rotation number of a rounding using turns of magnitude less than \(\pi\).
A polygon is \emph{uniform} when every turn is left or every turn is right.

The \emph{generic locus} is the set of polygons satisfying two
conditions: no three vertices are collinear (G1), and no point lies in
the relative interiors of three distinct edges (G2). On it, vertices are distinct, turns are nonzero,
nonincident crossings are transverse double points, and no vertex lies on
a nonincident edge. We first prove the theorem on the \emph{auxiliary
generic domain}: the subset of the generic locus on which, in addition,
every determinant appearing in the finite calculations below is nonzero,
including those of the finitely many smoothing components, whose vertices
include crossing points, and those of the auxiliary polygons of
Section~\ref{ep:section}, away from the explicitly analyzed endpoint
events. No value of \(\Theta(0)\) is used. This auxiliary domain is open
and dense in every chamber of the generic locus; the extension of the
theorem to the whole generic locus is carried out after the proof, in
Section~\ref{sec:completion}, by chamber constancy of both sides.

The amplitude \(A(P)\) is the finite marked-tree recursion of
GLSSW~\cite{GLSSW}: the marked root carries \(V^-\), ordinary internal
vertices carry \(-V^+\), the marked unary term is retained, and ordinary
unary insertions are excluded. We use GLSSW's root independence, or
cyclic invariance, as an input~\cite{GLSSW}: the tree coefficient does
not depend on the choice of root edge. Every polygon to which we apply
it below, namely smoothing components, caps closed by a straight segment,
and truncated or augmented polygons, is a generic polygon for generic
auxiliary data, which the auxiliary generic domain provides. We restate
the finite recursion and prove the auxiliary-current identities below.
We do not assume the equality to any knot expression.

\subsection{The positive diagram and its corner}

At every crossing choose the overpassing strand so that
\begin{equation}
 \det(t_{\mathrm{over}},t_{\mathrm{under}})>0.
 \label{eq:positive-diagram}
\end{equation}
This defines the positive diagram \(K_Q\) of an oriented subpolygon \(Q\).
It requires no contact lift. We retain this convention throughout; no
additional mirror operation is made in the proof or in the Wilson-loop
dictionary. The contact convention of the current letter can therefore be
kept as chosen there.

Our HOMFLY normalization is
\begin{equation}
 H_{\bigcirc}=1,
 \qquad aH_{L_+}-a^{-1}H_{L_-}=zH_{L_0},
 \qquad H_{L\sqcup\bigcirc}=\frac{a-a^{-1}}{z}H_L.
 \label{eq:homfly-conventions}
\end{equation}
For a uniform \(Q\), let \(m_Q\) count its retained self-crossings and put
\begin{equation}
 R_Q=|r(Q)|,
 \qquad d_Q=1-m_Q-R_Q,
 \qquad c(Q)=\coeff{a^{d_Q}z^0}H_{K_Q}(a,z).
 \label{eq:corner-definition}
\end{equation}
This is a \emph{prescribed} coefficient. Its definition does not replace
\(d_Q\) by the lowest occupied degree. No saturation or nonvanishing is
assumed. The proof below establishes the needed floor for the \(z^0\)
row of uniform polygons directly from descending skein expansion; it does
not require a bound on the entire HOMFLY polynomial.

\subsection{Separating decompositions}

Represent each crossing by a chord joining its two visits on the oriented
parameter circle of \(P\). Let \(\Dcal(P)\) consist of sets \(S\) of
crossings whose chords do not alternate pairwise. Equivalently, smoothing
every crossing in \(S\) separates the curve into exactly \(|S|+1\)
components; no smoothing rejoins two components. Denote the resulting
oriented polygons by \(P_S\). When evaluating an individual component we
retain its self-crossings and ignore its crossings with other components.
The latter convention will also follow from the lowest link row in the
framed skein calculation.

Let \(\Dcal_{\mathrm u}(P)\) be the subsets for which every member of
\(P_S\) is uniform. The state sum to be proved is
\begin{equation}
 C(P)=(-1)^{\ell(P)}
 \sum_{S\in\Dcal_{\mathrm u}(P)}(-1)^{|S|}
 \prod_{Q\in P_S}c(Q).
 \label{eq:state-sum}
\end{equation}
The empty set is included whenever the original polygon is uniform.
Each oriented smoothing replaces a crossing by one left and one right
corner. Original corners survive, so
\begin{equation}
 \sum_{Q\in P_S}\ell(Q)=\ell(P)+|S|.
 \label{eq:left-count}
\end{equation}

\begin{theorem}[Direct sewing identity]
\label{thm:main}
For every generic oriented polygon in the domain above, the GLSSW
amplitude satisfies \(A(P)=C(P)\), with conventions
\eqref{eq:positive-diagram}--\eqref{eq:state-sum}.
\end{theorem}

\paragraph{Relation to the letter.}
\begingroup\raggedright
The translation to the current v7 source is
\vseven{threebrackets}, \vseven{amplitudevertex}, and \vseven{amplitude}
for the tree recursion; \vseven{skein} and \vseven{skein2} for
\eqref{eq:homfly-conventions}; \vseven{factor} for
\eqref{eq:corner-definition}; and \vseven{C} for
\eqref{eq:state-sum}. These are identification labels, not external
equation references needed to compile or read this note. Neither
\vseven{softuniqueness} nor the letter's knot-degree theorem is imported.
\par\endgroup

\paragraph{Meaning of ``direct''.}
The original polygon is kept fixed while auxiliary endpoints are moved
inside finite interval-current identities. Determinant signs therefore
appear in the calculation, but no comparison of the chamber jumps of
\(A\) and \(C\), no connectivity argument for polygon chambers, and no
wall-law uniqueness principle enters. The string interpretation in the
last section is logically subsequent to the theorem.

% ---- endpoint ----
\section{Finite currents and transport of a marked endpoint}
\label{ep:section}

This section derives the geometric cutting identities directly from the
finite GLSSW recursion.  The only structural property of the GLSSW amplitude
used in addition to that recursion is its cyclic, or root, independence.
In particular, no wall law for a polygon invariant and no comparison with a
knot invariant is an input.

We first work with point-generic vertices: no three vertices occurring in a
sign evaluation are collinear.  Auxiliary moving points are evaluated on
the two sides of their isolated exceptional positions.  A value at the
exceptional position always means an expression which is unchanged between
these sides; it does not mean substituting zero in a sign function.

\subsection{A finite change of the near data}
\label{ep:near-data}

Cut an oriented polygon at a root edge and list the complementary boundary
word, including both root endpoints, as $q_0,\ldots,q_N$. Use the Cartesian convention
\begin{equation}
 H(i,k,j)=\sgn\det(q_k-q_j,q_i-q_j),\qquad i<k<j.
 \label{ep:H}
\end{equation}
The geometric near sign is $D(a,b,c)=H(a,b,c)$, the negative of the
Cartesian turn sign at that face corner. For signs $D,H\in\{1,-1\}$,
the GLSSW ordinary and marked gates obey
\begin{equation}
 D\Theta(-HD)=\frac{D-H}{2},\qquad
 D\Theta(HD)=\frac{D+H}{2}.
 \label{ep:linear-gates}
\end{equation}
We will be free to replace the near signs by an auxiliary array
$D(a,b,c)$. This is an algebraic change of interval variables, without a
claim that it has already been realized as a physical gauge symmetry.
For a strictly upper triangular array $X$ set
\begin{equation}
 \Phi_{D,H}(X)_{ij}=
 \sum_{i=r_0<\cdots<r_m=j}
 \left(\prod_{\nu=1}^{m-1}
 \frac{D(r_{\nu-1},r_\nu,r_{\nu+1})-H(i,r_\nu,j)}{2}\right)
 \prod_{\nu=1}^{m}X_{r_{\nu-1},r_\nu}.
 \label{ep:Phi}
\end{equation}
Empty products equal one.  Let $E_{i,i+1}=1$ and $E_{ij}=0$ for
$j>i+1$.  The ordinary current $b$ and the marked amplitude are precisely
\begin{equation}
 \Phi_{D,H}(b)=E,\qquad
 A(q_0,\ldots,q_N)=\Phi_{D,-H}(b)_{0N}.
 \label{ep:tree-recursion}
\end{equation}
The minus sign on each ordinary GLSSW vertex is incorporated by moving the
nonunary terms of the first equation to its right-hand side.  The marked
root includes its one-branch term.  Since that term has coefficient one,
the first equation has a unique solution, obtained by increasing interval
length.

\begin{lemma}[Near-data independence]
\label{ep:gauge}
The marked output in \eqref{ep:tree-recursion} is independent of $D$ for
fixed $H$.
\end{lemma}
\begin{proof}
Write $\mathcal F_H=\Phi_{0,H}$ and $\mathcal G_D=\Phi_{D,0}$.
A term in $\mathcal F_H\circ\mathcal G_D$ is specified by a final cut
list and a subset of its internal cuts: the subset records the outer
cuts.  An outer cut contributes $-H/2$, and an inner cut contributes
$D/2$.  The immediate neighbours used by an inner near entry are its
neighbours in the final list.  Conversely the list and subset uniquely
recover the outer and inner compositions.  Summing over subsets therefore
gives
\begin{equation}
 \Phi_{D,H}=\mathcal F_H\circ\mathcal G_D.
 \label{ep:factorization}
\end{equation}
All these maps have finite triangular inverses.  Solving the ordinary
equation first gives $\mathcal G_D(b)=\mathcal F_H^{-1}(E)$.
Substitution in the marked equation gives
\begin{equation}
 A=\bigl[\mathcal F_{-H}\mathcal F_H^{-1}(E)\bigr]_{0N},
 \label{ep:near-free}
\end{equation}
which contains no $D$.
\end{proof}

\subsection{An exact short-edge identity}
\label{ep:short-edge-section}

\begin{lemma}[Short-edge insertion]
\label{ep:short-edge}
Insert $B$ immediately after a vertex $A$ of a point-generic core polygon,
in a sufficiently short fixed direction.  Write $P,C$ for the predecessor
and successor of $A$ in the core. Define the short-edge sign explicitly by
$t(X)=\sgn\det(A-B,X-B)$, and put
\begin{equation}
 \alpha=t(P),\quad \beta=t(C),\quad
 \kappa=\frac{\alpha+\beta}{2}.
 \label{ep:kappa}
\end{equation}
These are the two attachment-turn signs in the GLSSW convention.  Then
\begin{equation}
 A(\text{polygon with }B)=\kappa A(\text{core}).
 \label{ep:duplication}
\end{equation}
\end{lemma}
\begin{proof}
Sufficient shortness means that every far triple not containing both
$A,B$ has its collapsed core sign.  This follows by continuity of finitely
many nonzero determinants.  Use Lemma~\ref{ep:gauge} to choose core near
entries $D^0(X,A,Y)=t(X)$.  Pull them back elsewhere, and set the new
entries $D^1(X,A,B)=t(X)$ and $D^1(A,B,Y)=t(Y)$.

The proposed ordinary parent current has the following exhaustive
description in terms of the core current $b^0$.  Intervals not containing
both labels have their collapsed value; $b^1_{AB}=1$; and
\begin{align}
 b^1_{AY}&=\frac{\beta-t(Y)}2 b^0_{AY}, &&Y>B,\notag\\
 b^1_{XB}&=\frac{\alpha-t(X)}2 b^0_{XA}, &&X<A,\notag\\
 b^1_{XY}&=\kappa b^0_{XY}, &&X<A<B<Y.
 \label{ep:duplicate-current-table}
\end{align}
To verify the first endpoint formula, group parent cut lists according to
their collapsed list and let $X$ be its first cut after $A$.  The two
presentations, with the short edge inside the first branch or as a separate
leaf, have multipliers
\begin{equation}
 \frac{\beta-t(X)}2+\frac{t(X)+\eta t(Y)}2
 =\frac{\beta+\eta t(Y)}2,
 \label{ep:endpoint-two-lifts}
\end{equation}
where $\eta=-1$ is ordinary and $\eta=+1$ is marked.  For the ordinary
equation this multiplies the core source $E$: a longer interval has source
zero, and the one-edge case has $Y=C$, so $t(Y)=\beta$.  The other endpoint
formula follows by reversing this bookkeeping, with $\alpha$ in place of
$\beta$.

For a strictly spanning interval, a collapsed list without a cut at $A$
has one presentation and acquires $\kappa$.  If it cuts at $A$, let $L,R$
be the neighbouring cuts and put $a=t(L)$, $b=t(R)$, and
$h=H(\text{interval start},A,\text{interval end})$.  The three
presentations cut only at $B$, only at $A$, or at both.  Their extra child
multipliers are respectively $(\alpha-a)/2$, $(\beta-b)/2$, and $1$;
their replacing gates are $(a+\eta h)/2$, $(a+\eta h)/2$, and
$(a+\eta h)(b+\eta h)/4$.  Subtracting $\kappa$ times the core gate
leaves
\begin{align}
 -\frac{(a+b)(a+\eta h)}4
 +\frac{(a+\eta h)(b+\eta h)}4
 &=\frac{h^2-a^2}{4}\notag\\
 &=0.
 \label{ep:three-lifts}
\end{align}
The last step uses $a,h\in\{1,-1\}$.  This verifies every coordinate of
the ordinary recursion and hence its unique solution.  The same grouping
with $\eta=+1$ proves the marked formula.  Root independence allows a root
different from the short edge throughout the calculation.
\end{proof}

\begin{lemma}[Exposed truncation]
\label{ep:truncation}
Replace an exposed Cartesian left corner by two nearby points on its
incident edges, joined by a short supporting edge.  The GLSSW amplitude
changes sign.  The same is true of any separating-smoothing alternating
sum whose component weights are $A$.  In particular, multiplication by
$(-1)^{\#\text{vertices}}$ removes the change.
\end{lemma}
\begin{proof}
Insert a short edge inside the exposed wedge, with both attachment-turn
signs equal to $-1$.  Lemma~\ref{ep:short-edge} gives the multiplier
$-1$.  The insertion can be placed on the two incident lines without
changing any nonzero far signs.  Choose its neighbourhood disjoint from
all crossings and other edges.  Every separating support then survives
unchanged and exactly one of its components contains the replaced corner.
Apply the same short-edge identity to that component.  One vertex has
been added, which supplies the second minus sign in the normalized sum.
\end{proof}

\subsection{The endpoint cut-list calculation}
\label{ep:endpoint-calculation}

Let $q_1,\ldots,q_m$ be a fixed open arc.  In this subsection take
middle-only near data $D(a,b,c)=d_b$, with arbitrary fixed signs $d_b$.
For a moving point $p$, define the unit upper triangular matrix
\begin{equation}
 F(p)_{ai}=b(p,q_a,\ldots,q_i),\qquad F(p)_{aa}=1.
 \label{ep:F}
\end{equation}
The value on the diagonal is the one-edge ordinary current.  All matrices
below are finite.

Suppose $p$ crosses the line through $q_k,q_j$ at an isolated event, and
write $q_k=p+\tau(q_j-p)$ there.  Write
$\Delta H(p,q_k,q_j)=2\delta$.  The ordinary recurrence is linear in
the first-endpoint currents.  At its directly changed gate, split the cut
list into the prefix ending at $k$ and the suffix beginning at $k$.
The identities of determinants at the collinear position give
\begin{equation}
 H(p,q_i,q_k)=\sgn(\tau)H(p,q_i,q_j),\qquad
 H(q_k,q_i,q_j)=\sgn(1-\tau)H(p,q_i,q_j).
 \label{ep:collinear-far-identities}
\end{equation}
Thus the prefix sum is the ordinary source $E$ if $\tau>0$ and the
marked amplitude if $\tau<0$.  The suffix sum is the ordinary source
$E$ if $\tau<1$ and the marked amplitude if $\tau>1$.  An ordinary
source can be nonzero only on a one-edge interval.  It follows that the
direct source, and its propagated endpoint change, are
\begin{equation}
 \Delta F=
 \begin{cases}
 \delta\,a^{(k)}e_{k+1}^{T}F,&\tau<0,\ j=k+1,\\
 \delta\,e_k e_{k+1}^{T}F,&0<\tau<1,\ j=k+1,\\
 \delta\,e_k A(q_k,\ldots,q_j)e_j^{T}F,&\tau>1,
 \end{cases}
 \label{ep:three-endpoint-events}
\end{equation}
All other cases vanish.  Here $a^{(k)}_a=A(p,q_a,\ldots,q_k)$ for
$a\leq k$, with $a^{(k)}_k=1$, and its remaining entries vanish.
For completeness, the propagation factor $e_j^TF$ comes from grouping a
larger endpoint cut list by the branch carrying the changed source.  The
part before that branch has already been included in the source.  The
remaining cuts are exactly the ordinary recurrence on
$p,q_j,\ldots,q_i$, with unit first branch.  Middle-only near data
make this immediate: the near factor at the first cut does not depend on
the identity of the first branch, whose predecessor is always $p$.  This
gives $F_{ji}$ by triangular uniqueness.  The minus sign of the ordinary top recurrence cancels the
minus sign in the gate change $\Delta[(d-H)/2]=-\delta$.

\subsection{A cap identity along a fixed line}
\label{ep:line-caps}

Take the reference line to be $y=0$ and assume no original vertex lies on
it.  Set
\begin{equation}
 h_i=\sgn y_i,\quad S=\operatorname{diag}(h_i),\quad
 B_{ij}=b(q_i,\ldots,q_j)\ (i<j),\quad B_{ii}=0,\quad R=(I-BS)^{-1},
 \label{ep:line-data}
\end{equation}
using the near choice $D(a,b,c)=h_b$.  For $p(t)=(t,0)$ let
$F(t)$ be \eqref{ep:F}, and define
\begin{equation}
 G(t)_{ib}=b(q_i,\ldots,q_b,p(t)),\quad G(t)_{ii}=1,
 \qquad W(t)=F(t)SRG(t).
 \label{ep:G-W}
\end{equation}

\begin{lemma}[Short-edge pairing]
\label{ep:wronskian}
For $a<b$,
\begin{equation}
 W(t)_{ab}=\frac{h_a+h_b}{2}\,A(p(t),q_a,\ldots,q_b).
 \label{ep:W-entry}
\end{equation}
The diagonal has $W_{aa}=h_a$.  In particular $SW=WS$.
\end{lemma}
\begin{proof}
Insert $p'=p-\epsilon(1,0)$ immediately after $p$ and root on the short
edge $p\to p'$.  On the complementary word
$p',q_a,\ldots,q_b,p$, its far signs are $h_i$.  Its ordinary top
gates therefore vanish, while its marked top gates are $h_i$.  The term
of the marked expansion with no internal cut is the ordinary top current
of the word, which vanishes with those gates.  Every other term begins
with a first branch $F(p')_{ai}$, equal to $F(p)_{ai}$ for small
$\epsilon$, and ends with a last branch $G(t)_{jb}$.  Expanding
over the first and last internal cuts therefore gives $F S R G$:
each additional internal branch contributes $BS$.  Lemma~\ref{ep:short-edge}
and root independence evaluate the same polygon as the right-hand side
of \eqref{ep:W-entry}.  For $a=b$ this is the elementary short triangle;
equivalently one may assign the formal two-edge core value one.
If $h_a\ne h_b$ the right-hand side vanishes.  This proves commutation
with the diagonal matrix $S$ entry by entry.
\end{proof}

Put
\begin{equation}
 J(t)=F(t)^{-1}SF(t),\qquad U(t)=SRG(t).
 \label{ep:J-U}
\end{equation}
The second identity below is the definition of $J$ rewritten. Since
$W=FU$ and $F$ is invertible, the commutation $SW=WS$ gives
$FJU=SFU=SW=WS=FUS$, hence the first:
\begin{equation}
 J U=U S,\qquad F J=S F.
 \label{ep:intertwining}
\end{equation}
The two virtual types in \eqref{ep:three-endpoint-events} have $h_k=h_j$:
$q_k,q_j$ lie on the same side of the line.  Their elementary left
multiplier of $F$ therefore commutes with $S$ and leaves $J$ unchanged.

Let $c$ instead be an actual crossing of the line by the oriented edge
$q_k\to q_{k+1}$.  Its direction is
$s_c=h_{k+1}=-h_k$.  Moving $p$ from left to right changes the critical
far sign by $2s_c$.  Thus \eqref{ep:three-endpoint-events} gives
$\Delta F=s_c a^{(k)}e_{k+1}^{T}F$.  Define
\begin{equation}
 u_c=U(c)e_k,\qquad v_c=e_{k+1}^{T}F(c).
 \label{ep:uv}
\end{equation}
These expressions omit one or the other endpoint of the critical edge,
so their values are the same on the two sides of the event.
Column $k$ of \eqref{ep:W-entry} says
$F u_c=(S-s_c I)a^{(k)}/2$.  The square of
$a^{(k)}e_{k+1}^T$ is zero, by its triangular support.  Write $N=a^{(k)}e_{k+1}^{T}$, so that
$F_{\rm after}=(I+s_cN)F_{\rm before}$ and $(I+s_cN)^{-1}=I-s_cN$.
Conjugating $S$ by this elementary multiplier gives
\begin{equation}
 \Delta J
 =F_{\rm before}^{-1}\bigl[s_c(SN-NS)-s_c^2NSN\bigr]F_{\rm before}
 =s_cF_{\rm before}^{-1}(S-s_cI)a^{(k)}e_{k+1}^{T}F_{\rm before}
 =2s_c u_cv_c.
 \label{ep:J-jump}
\end{equation}
Here $NS=s_cN$ because $e_{k+1}^{T}S=h_{k+1}e_{k+1}^{T}$ and
$h_{k+1}=s_c$; there is no quadratic term because
$NSN=a^{(k)}(e_{k+1}^{T}Sa^{(k)})e_{k+1}^{T}=0$, as $a^{(k)}$ is
supported on indices at most $k$; and the last step is
$F u_c=(S-s_cI)a^{(k)}/2$ together with $v_c=e_{k+1}^{T}F$.  The intertwining identities give, on either side,
\begin{equation}
 v_cu_c=0,\qquad J(c)u_c=-s_cu_c,\qquad v_cJ(c)=s_cv_c.
 \label{ep:rank-one-identities}
\end{equation}

Order the actual crossings by their horizontal coordinate.  For $x_c<x_d$
sum \eqref{ep:J-jump} between them and sandwich the result between
$v_c$ and $u_d$.  Virtual events contribute zero.  The jumps at $c$ and at $d$ themselves
drop out, since they are sandwiched into $v_cu_c=0$ and $v_du_d=0$ by
\eqref{ep:rank-one-identities}.  The resulting identity is
\begin{equation}
 (s_c+s_d)v_cu_d+
 2\sum_{x_c<x_e<x_d}s_e(v_cu_e)(v_eu_d)=0.
 \label{ep:pairwise-cap}
\end{equation}
Indeed the left and right endpoint values of $J$ act respectively by
$s_c$ on the row and $-s_d$ on the column, which accounts for the sign.

When $c$ occurs before $d$ on the original arc, the cap consists of that
subarc followed by the straight segment $d\to c$.  Since $x_c<x_d$, this
closing segment runs leftward, and
$\operatorname{sgn}\det(q_i-d,\,c-d)=\operatorname{sgn}y_i=h_i$: all the
far signs of the cap at its horizontal closing edge are $h_i$.  Its ordinary top current vanishes
and its marked top expansion is exactly $v_cSRG(d)e_{k_d}=v_cu_d$.
Thus $v_cu_d=A(\text{cap }c\to d)$.  If the arc order is reversed,
the pairing vanishes by triangular support.  Let $X$ be the strictly
upper triangular matrix of these pairings, indexed by horizontal order,
and let $S_C=\operatorname{diag}(s_c)$.  Equation~\eqref{ep:pairwise-cap}
is precisely
\begin{equation}
 S_CX+XS_C+2XS_CX=0.
 \label{ep:cap-matrix}
\end{equation}
Define $K=(I+X)^{-1}$.  Multiplying \eqref{ep:cap-matrix} on both sides
by $K$ gives
\begin{equation}
 KS_C+S_CK=2S_C.
 \label{ep:K}
\end{equation}
One can check the multiplication without commuting any matrices:
$(I+X)S_C+S_C(I+X)-2(I+X)S_C(I+X)$ is the negative of the left-hand
side of \eqref{ep:cap-matrix}.  Therefore $K_{cd}=0$ for $c\ne d$
with $s_c=s_d$.

\begin{theorem}[Mixed-edge recursion]
\label{ep:mixed-edge}
Let $e$ be an edge of $P$ whose two endpoint turns are opposite.  Then
\begin{equation}
 A(P)=\sum_{\substack{\varnothing\ne T\subset\{\text{crossings on }e\}\\
                         T\ \mathrm{separating}}}
 (-1)^{|T|+1}\prod_{Q\in P_T}A(Q).
 \label{ep:mixed-edge-formula}
\end{equation}
In particular the amplitude vanishes if such an edge is uncrossed.
\end{theorem}
\begin{proof}
Choose coordinates so that $e$ closes the remaining arc from its right
endpoint $v$ to its left endpoint $u$.  Extend the remaining arc slightly
backwards past $u$ along its outgoing edge and forwards past $v$ along its
incoming edge.  Omit $u,v$ as vertices of the extended arc: they are now
two actual crossings of the reference line.  The edges at $u$ and at $v$
other than $e$ are not horizontal, since the turns there are nonzero, and
by (G1) no other vertex lies on the line through $e$; so the extended arc
has no vertex on the reference line, $u$ and $v$ are transverse crossings,
and no flat external leg is introduced.  The cap from $u$ to $v$ is
exactly $P$.

If the first and last vertices of the return arc have height signs
$h_{\rm first},h_{\rm last}$, the crossing directions at the two ends
are $s_u=h_{\rm first}$ and $s_v=-h_{\rm last}$.  The Cartesian turn
signs of $P$ at these ends are respectively $-h_{\rm first}$ and
$-h_{\rm last}$.  Opposite turns therefore give $s_u=s_v$, so
$K_{uv}=0$ by \eqref{ep:K}.

Expand the finite inverse as $K=I-X+X^2-\cdots$.  Its $uv$ entry starts
with $-A(P)$.  A later term is a chain of intermediate crossings
increasing both in horizontal order and in arc order; its factors are the
successive cap amplitudes.  Such chains are exactly the separating
smoothing supports on $e$: the visits along the closing edge have the
reverse horizontal order, so increasing arc order is equivalent to
noninterlacement of their Gauss chords.  The components of $P_T$ are
these caps: the oriented smoothing at a selected crossing $c_i$ sends the
arc arriving at $c_i$ onto $e$, toward $u$, and the strand arriving along
$e$ onto the outgoing arc, so each component is the subarc between two
consecutive selected crossings closed along $e$, traversed in the
orientation of $P$.  A chain of $r$ intermediate cuts
contributes $(-1)^{r+1}$ times its $r+1$ cap factors.  Moving the first
term to the other side gives \eqref{ep:mixed-edge-formula}.
\end{proof}

Coincident secant events, or two crossings of the reference line at the
same horizontal coordinate, are excluded on the auxiliary generic domain
of Section~\ref{sec:statement}: a coincidence on the segment $e$ itself
would be a triple point, excluded by (G2), and the remaining ones are
vanishing auxiliary determinants.  Both sides of
\eqref{ep:mixed-edge-formula} depend only on the chirotopes of the
finitely many caps, and the identity extends to the whole generic locus
by the continuation argument of Section~\ref{sec:completion}.

\begin{remark}[Harmless far-sign degeneracies]
For a closed polygon satisfying (G1), an accidental collinear triple of
the auxiliary calculation either contains no boundary edge or has its
remote vertex outside the boundary segment. Root the polygon on an edge
outside the triple and choose middle-only near data, which
Lemma~\ref{ep:gauge} permits; the triple then enters only as a far sign.
In the cut-list calculation leading to \eqref{ep:three-endpoint-events},
the prefix or suffix factor that is an ordinary source is then evaluated
on a nontrivial interval and vanishes. The full marked amplitude
consequently has the same value on both generic sides. This is the
chamber constancy of $A$ seen from the recursion; individual auxiliary
currents need not share it. Section~\ref{sec:completion} does not rely on
this remark: the extension of the theorem to the whole generic locus is
by chamber constancy of both sides. No prescription for $\Theta(0)$ is
being imposed.
\end{remark}

\subsection{Transport on the actual boundary}
\label{ep:boundary-section}

For the uniform-vertex argument a different auxiliary choice is useful:
take $D=-1$ everywhere, as Lemma~\ref{ep:gauge} permits.  For an open
word define
\begin{equation}
 M(p,q_1,\ldots,q_n)
 =\bigl[B(I+B)^{-1}\bigr]_{0n},
 \label{ep:M}
\end{equation}
where index $0$ denotes $p$, and $B_{ij}$ is the ordinary current on that
subword.  Equivalently $M$ is the sum over all interval cut lists with
weight $(-1)^{\#\text{cuts}}$.  If all the other vertices lie strictly
to the left of the directed closing edge $q_n\to p$, then its far signs
$H(0,k,n)=\sgn\det(q_k-q_n,\,p-q_n)$ are all $-1$.
The marked root gates $(D+H)/2$ are then all $-1$, and consequently $M=A$.

\begin{lemma}[Boundary crossing transport]
\label{ep:boundary-jump}
Move $p$ forward along the actual first edge toward $q_1$.  Suppose that
edge meets $q_k\to q_{k+1}$ at $c$, and put
$\delta_c=\sgn\det(q_1-p,q_{k+1}-q_k)$.  The two one-sided currents
satisfy
\begin{equation}
 \Delta M=
 \delta_c A(c,q_1,\ldots,q_k)
 M(c,q_{k+1},\ldots,q_n).
 \label{ep:boundary-jump-formula}
\end{equation}
Here $\Delta$ is the value after passage minus the value before passage.
For the canonical positive lift, $\delta_c=-1$ at a first-under encounter
and $+1$ at a first-over encounter.
\end{lemma}
\begin{proof}
Apply \eqref{ep:three-endpoint-events}, whose hypothesis of middle-only
near data the constant choice $D=-1$ satisfies, to the first row, $a=1$;
the events along the first edge are isolated by genericity.
A virtual event must have $k=a$, hence $p,q_1,q_j$ collinear.
On the interior of the first physical edge this is impossible by point
genericity: its only meeting with that secant line is $p=q_1$.
The surviving event is the crossing of an actual edge and has
\begin{equation}
 \Delta b(p,q_1,\ldots,q_j)
 =\delta_c A(c,q_1,\ldots,q_k)
 b(c,q_{k+1},\ldots,q_j).
 \label{ep:boundary-b}
\end{equation}
Write the first row of $B$ as $f$ and the fixed suffix block as $B_*$.
Block multiplication in \eqref{ep:M} gives its first row as
$f(I+B_*)^{-1}$.  Right-multiplying \eqref{ep:boundary-b} by that fixed
inverse converts its parent $b$ to the parent $M$, proving
\eqref{ep:boundary-jump-formula}.  Finally the positive-lift convention is
$\det(t_{\rm over},t_{\rm under})>0$; inserting the tangent of the
moving first strand gives exactly the two signs asserted.
\end{proof}

\begin{lemma}[Turn projector]
\label{ep:turn}
When the first edge shrinks to its terminal vertex $q_1$, the current on
the remaining open word is multiplied by $-1$ at a Cartesian left turn
and by $0$ at a Cartesian right turn.  The one-edge current by itself has
value one.
\end{lemma}
\begin{proof}
Let $p$ approach $q_1$ along the incoming edge and put
$\sigma=H(p,q_1,q_2)$.  Take the final segment sufficiently short so that
all far signs involving $p$ but not both $p,q_1$ agree with their collapsed
signs; the far signs involving both $p$ and $q_1$ are constant along the
approach, since $p$ moves on the line through $q_1$ and no other vertex
lies on it.  Group the ordinary cut lists according to whether they cut at
$q_1$.  The same two-presentation calculation as
\eqref{ep:endpoint-two-lifts}, now with constant near data $-1$, gives
\begin{equation}
 b(p,q_1,\ldots,q_j)=\eta\,b(q_1,\ldots,q_j),\qquad
 \eta=\frac{1+\sigma}{2},\qquad j>1,
 \label{ep:short-first-current}
\end{equation}
while $b(p,q_1)=1$.  For $j=2$ this is the direct computation
$b(p,q_1,q_2)=-(-1-\sigma)/2=\eta$.  For $j>2$ it follows by induction on
$j$: group the cut lists of $(p,q_1,\ldots,q_j)$ by whether they cut at
$q_1$.  Those that do not have the gates of the cut lists of
$(q_1,\ldots,q_j)$, with the first branch replaced by
$b(p,q_1,\ldots,q_{r_1})=\eta\,b(q_1,\ldots,q_{r_1})$, and their sum is
$\eta$ times the sum defining $-b(q_1,\ldots,q_j)$.  Those that do carry
the gate $(-1-H(p,q_1,q_j))/2$ times the full ordinary recursion of
$(q_1,\ldots,q_j)$ including its one-part term, which is the source
$E=0$ by \eqref{ep:tree-recursion}.

Now use the cut-list expansion of $M$.  Lists without a cut at $q_1$
supply $\eta M(q_1,\ldots,q_j)$.  Lists with that cut supply
$-M(q_1,\ldots,q_j)$, since the initial one-edge current is one and the
new cut has weight $-1$.  Their sum is
\begin{equation}
 M(p,q_1,\ldots,q_j)
 =(\eta-1)M(q_1,\ldots,q_j).
 \label{ep:turn-formula}
\end{equation}
The sign convention \eqref{ep:H} makes $\sigma=-1$ at a Cartesian left
turn and $+1$ at a right turn.  Hence $\eta-1$ is respectively $-1$ and
$0$.
\end{proof}

% ---- primitive ----
\section{Separating-smoothing inversion and nonuniform support}
\label{sec:primitive}

For any scalar function \(X\) of closed polygons define its separating
primitive by
\begin{equation}
 \Gamma_X(P)=\sum_{T\in\Dcal(P)}(-1)^{|T|}
                    \prod_{Q\in P_T}X(Q).
 \label{eq:primitive}
\end{equation}
The term ``primitive'' here refers to subtraction of separating sewings.
It is not a claim about a connected Wilson cumulant or a one-particle
irreducible Green function.

\begin{lemma}[Finite inversion]
\label{lem:inversion}
For every \(P\),
\begin{equation}
 X(P)=\sum_{S\in\Dcal(P)}\prod_{Q\in P_S}\Gamma_X(Q).
 \label{eq:inversion}
\end{equation}
\end{lemma}
\begin{proof}
Expand the right side using \eqref{eq:primitive}. A term first cuts a
support \(S\), then chooses independent separating cuts inside each
component. Their union is a final separating support \(U\). Conversely,
for a fixed \(U\), every subset \(S\subseteq U\) gives exactly one such
two-stage refinement: nonalternation is inherited by subsets and by the
components after smoothing. The coefficient of
\(\prod_{Q\in P_U}X(Q)\) is therefore
\begin{equation}
 \sum_{S\subseteq U}(-1)^{|U|-|S|}.
 \label{eq:boolean-cancellation}
\end{equation}
It is zero for nonempty \(U\) and one for \(U=\varnothing\). The surviving
term is \(X(P)\).
\end{proof}

\begin{proposition}[Support of the GLSSW primitive]
\label{prop:nonuniform}
If \(P\) is not uniform, then \(\Gamma_A(P)=0\).
\end{proposition}
\begin{proof}
Some edge \(e\) has opposite turn signs at its two endpoints. Group
\eqref{eq:primitive} by the subset \(U\) of selected crossings that are
not on \(e\). In \(P_U\), the whole edge \(e\) belongs to one component
\(R\); its two original endpoint turns still have opposite signs.
Every other component is a spectator.

The remaining admissible cuts are exactly separating subsets of the
retained self-crossings of \(R\) on \(e\). Crossings between \(R\) and a
spectator cannot be chosen, since they would rejoin components. Thus the
sum multiplying the spectator factors is
\begin{equation}
 A(R)+
 \sum_{\substack{\varnothing\ne T\in\Dcal(R)\\T\subset e}}
       (-1)^{|T|}\prod_{Q\in R_T}A(Q).
 \label{eq:mixed-edge-cancellation}
\end{equation}
Lemma~\ref{ep:mixed-edge}, proved from the interval recursion above, makes
\eqref{eq:mixed-edge-cancellation} zero. This holds for each fixed \(U\),
so the entire primitive vanishes.
\end{proof}

By inversion, this already proves a state sum over uniform pieces with
weights \((-1)^{\ell(Q)}\Gamma_A(Q)\). The remaining question is to
identify those weights. The next section does so without using the
state-sum equality.


% ---- uniform ----
\section{The uniform primitive is the HOMFLY corner}
\label{un:section}

This section identifies the uniform weights left undetermined by the
crossing recursion.  The argument compares two finite recursions: a
boundary traversal of a marked GLSSW current, and the descending skein
calculation of a framed HOMFLY coefficient.  No polygon wall law, equality
$A=C$, or transverse-knot degree bound is used.

Write
\begin{equation}
 \nu_A(Q):=(-1)^{\ell(Q)}\Gamma_A(Q).
 \label{un:nu}
\end{equation}
We shall prove $\nu_A(Q)=c(Q)$ for uniformly turning $Q$.  We first orient
$Q$ so that all its corners turn left; then $\ell(Q)$ is its number of
vertices.  Orientation reversal is treated at the end.

\subsection{An open component and its genuine smoothing supports}
\label{un:open}

For a polygonal arc $\gamma=(p,q_1,\ldots,q_N)$, let $B$ denote the
strictly upper triangular matrix of ordinary interval currents in the
auxiliary choice $D=-1$.  Define
\begin{equation}
 M:=B(I+B)^{-1}=B-B^2+B^3-\cdots,
 \qquad M(\gamma):=M_{0N}.
 \label{un:M}
\end{equation}
The sum is finite.  Only the actual edges of $\gamma$ are allowed to
participate in a smoothing.  In particular, the auxiliary straight
segment from $q_N$ to $p$, implicit in an interval current, supplies
\emph{no} smoothing sites.

Let $\mathcal D(\gamma)$ be the sets of pairwise nonalternating
self-crossing chords of the real arc.  Smoothing $T\in\mathcal D(\gamma)$
leaves one open component $\gamma_T$ and $|T|$ closed components.  Set
\begin{equation}
 \Gamma_M(\gamma):=
 \sum_{T\in\mathcal D(\gamma)}(-1)^{|T|}
 M(\gamma_T)\prod_{Q\text{ closed in }\gamma_T}A(Q).
 \label{un:GammaM}
\end{equation}
The notation for the product means the closed components of the entire
smoothed arc, not components of its open member.

There is a useful genericity observation.  A component of a separating
smoothing contains at most one connected segment of any given original
edge.  To see this, root the Gauss circle immediately before that edge.
All selected visits on it occur consecutively, say $a_1<\cdots<a_t$.
Their other visits must occur in the opposite order: otherwise two
selected chords alternate.  The resulting nested cuts place successive
segments of the edge in different components.  Thus a component cannot
acquire three collinear vertices merely from several cuts on one edge.
All remaining accidental point-triple incidences can be excluded by
genericity.  The short-edge limits below are taken one-sided; they are
not evaluations of a GLSSW sign at zero.

Suppose that the closing edge $q_Np$ is an exposed supporting edge with
the rest of the polygon on its left.  Its far signs are all $-1$.
Consequently the marked-root expansion is precisely
$B-B^2+\cdots$, and $M(\gamma)=A(Q)$.  The same supporting edge survives in
the open component of every $T$.  Hence, term by term,
\begin{equation}
 \Gamma_M(\gamma)=\Gamma_A(Q).
 \label{un:initial}
\end{equation}
If $Q$ has no such edge, truncate a vertex of its convex hull, which is a
left corner with the polygon on its left, by a sufficiently small
supporting segment.  Lemma~\ref{ep:truncation} changes the amplitude
of the affected component by $-1$.  The smoothing supports are unchanged,
so it changes $\Gamma_A$ by $-1$ as well.  It adds one vertex, and
therefore preserves $\nu_A$.  On the knot side it changes neither the
lift, the retained crossing number, nor the rotation number.  We may
therefore start at an exposed edge without changing the desired identity.

\subsection{Moving the mark: the smoothing subtraction matters}
\label{un:transport}

Move $p$ forward along the first physical edge.  At a crossing $c$ with
the later edge $q_kq_{k+1}$, put
\begin{equation}
 \delta=\operatorname{sgn}\det(q_1-p,q_{k+1}-q_k),\qquad
 Q_c=(c,q_1,\ldots,q_k),\qquad
 \gamma_c=(c,q_{k+1},\ldots,q_N).
 \label{un:crossing-data}
\end{equation}
In the positive lift, $\delta=-1$ is a first-under encounter and
$\delta=+1$ is a first-over encounter.  A jump always means the value
after the mark has passed the crossing minus its value before.

\begin{lemma}[Boundary transport after smoothing subtraction]
\label{un:transport-lemma}
For the real-arc current \eqref{un:GammaM},
\begin{equation}
 \Delta\Gamma_M
   =\frac{\delta-1}{2}\,
       \Gamma_A(Q_c)\Gamma_M(\gamma_c).
 \label{un:primitive-jump}
\end{equation}
In particular, the current changes by minus the child--parent product
at a first-under crossing and does not change at a first-over crossing.
\end{lemma}

\begin{proof}
Separate the smoothing supports according to whether they contain $c$.
For supports not containing $c$, a selected chord that alternates with
the chord of $c$ puts its two visits in different components.  The open
component then has no self-crossing event at $c$; the passage of $p$
across the line of the crossed edge is at most a virtual secant event of
the open word, which by the proof of Lemma~\ref{ep:boundary-jump} never
reaches the moving first row, so its current has no jump.  For the
remaining supports, note that $c$ is the first crossing ahead of $p$ on
the first edge, and cut first at $c$.  The selected
chords lie independently in the resulting full closed child and open
parent.  Conversely, independent separating supports on those two
pieces uniquely reconstruct the original support.  This refinement
bijection preserves the sum of the support sizes.

Lemma~\ref{ep:boundary-jump}, applied to the open component in each
summand, therefore gives the contribution
\begin{equation}
 \delta\,\Gamma_A(Q_c)\Gamma_M(\gamma_c)
 \label{un:no-cut-contribution}
\end{equation}
from supports not containing $c$.

Supports containing $c$ exist immediately before the mark crosses $c$
and disappear immediately after.  Before the event their open parent
starts with the tiny edge $pc$, followed by $cq_{k+1}$; their extra
smoothing has weight $-1$.  At $c$ this tiny-edge parent turns left if
$\delta=+1$ and right if $\delta=-1$.  The short-edge projector of
Lemma~\ref{ep:turn} gives, term by term in its other smoothing supports,
\begin{equation}
 \lim_{p\to c^-}\Gamma_M(p,c,q_{k+1},\ldots,q_N)
   =-\frac{1+\delta}{2}\,\Gamma_M(\gamma_c).
 \label{un:tiny-parent}
\end{equation}
Its tiny initial edge has no crossings, so the projected corner always
belongs to the open component.  Removing these pre-event terms
contributes
\begin{equation}
 -\frac{1+\delta}{2}\,
       \Gamma_A(Q_c)\Gamma_M(\gamma_c)
 \label{un:cut-contribution}
\end{equation}
to the jump.  Adding \eqref{un:no-cut-contribution} and
\eqref{un:cut-contribution} gives \eqref{un:primitive-jump}.

This argument invokes the endpoint lemma only at actual boundary
events.  Its virtual-secant terms do not reach the moving first row.
Moreover, no term of \eqref{un:GammaM} selects a crossing of an auxiliary
closing segment.  Thus the calculation introduces neither fictitious
smoothing sites nor a hidden change of the closed knot.
\end{proof}

When the mark passes an original left corner, the crossing-free short
first edge projects $\Gamma_M$ to minus the current with that corner
removed.  Indeed, no chord visits the crossing-free first edge, so the
corner $q_1$ remains in the open component of every support, and
Lemma~\ref{ep:turn} applies termwise.  Normalize the current
by the number $v(\gamma)$ of vertices of its active word, including its
two endpoints:
\begin{equation}
 f(\gamma):=(-1)^{v(\gamma)}\Gamma_M(\gamma).
 \label{un:f}
\end{equation}
Left-corner passage then leaves $f$ unchanged.  At a first-under crossing
the child closing corner is left, so $Q_c$ is uniformly left-turning
whenever the unprocessed original corners are left.  Furthermore,
\begin{equation}
 v(Q_c)+v(\gamma_c)=v(\gamma)+1.
 \label{un:vertex-count}
\end{equation}
Multiplying \eqref{un:primitive-jump} by the vertex parity and using
\eqref{un:vertex-count} yields
\begin{equation}
 f_{\rm before}
   =f_{\rm after}-\nu_A(Q_c)f(\gamma_c)
       \quad\hbox{at a first-under crossing}.
 \label{un:f-recursion}
\end{equation}
The terminal arc has one edge, so $M=1$, there are no smoothing supports,
and $v=2$.  Its terminal value is therefore
\begin{equation}
 f_{\rm terminal}=1.
 \label{un:f-terminal}
\end{equation}

For a detached child we restart this construction at one of its own
exposed points.  This is permitted by the root independence of $A$ and
the cyclic definition of $\Gamma_A$.  Each detached child has strictly
fewer retained self-crossings than the piece from which it was detached,
and the recursion is finite: the pair (retained self-crossings, vertices
of the active word) decreases lexicographically at every event.

\subsection{The framed lowest HOMFLY row}
\label{un:framed}

We record the precise link normalization used in the comparison.  For
an oriented knot diagram $D$, including diagrams in which some crossings
have already been switched, let $w(D)$ be its signed writhe and $r(D)$
its signed plane rotation.  Define
\begin{equation}
 h_D(a):=[z^0]H_D(a,z),\qquad
 \mathcal B_D(a):=
 a^{w(D)+r(D)}(a-a^{-1})h_D(a),\qquad y:=a^2-1.
 \label{un:framed-definition}
\end{equation}
The writhe is diagram-dependent.  Its inclusion is essential even when
the underlying knot is an unknot.

\begin{lemma}[Lowest link row]
\label{un:lowest-link-row}
If $L=D_1\cup\cdots\cup D_k$ and
$\Lambda(L)=\sum_{i<j}\operatorname{lk}(D_i,D_j)$, then
\begin{equation}
 [z^{1-k}]H_L
 =a^{-2\Lambda(L)}(a-a^{-1})^{k-1}
     \prod_{i=1}^k h_{D_i}(a).
 \label{un:link-factor}
\end{equation}
Every $k$-component HOMFLY polynomial has no powers of $z$ below $1-k$.
\end{lemma}

\begin{proof}
The $z$ bound follows directly by descending skein induction.  A terminal
descending diagram is a split unlink and has lowest power $z^{1-k}$.
At a self-crossing the smoothing has $k+1$ components, so multiplication
by the skein factor $z$ restores the lower bound $1-k$.  At a crossing
between components the smoothing has $k-1$ components, so that same
factor gives the stronger lower bound $3-k$.  The induction is ordered
first by crossing number and then by the number of crossings encountered
first underneath; switching decreases the second number and smoothing
decreases the first.

At a mixed crossing, the smoothing term therefore cannot contribute to
$z^{1-k}$.  The skein relation gives
$[z^{1-k}]H_{L_+}=a^{-2}[z^{1-k}]H_{L_-}$.  Switching from negative to
positive raises $\Lambda$ by one.  Hence
$a^{2\Lambda}[z^{1-k}]H_L$ is invariant under mixed-crossing changes.
Make only such changes until the components are split.  The split-union
normalization gives $(a-a^{-1})^{k-1}\prod_i h_{D_i}$ there, proving
\eqref{un:link-factor}.  This is the lowest-row factorization in the
Lickorish--Millett normalization.
\end{proof}

Apply the skein relation at an untouched positive self-crossing of $D$.
Let $D^{\rm sw}$ be its switched diagram and let $D_L\cup D_R$ be its
oriented smoothing.  First extract the knot row:
\begin{equation}
 a h_D-a^{-1}h_{D^{\rm sw}}
   =a^{-2\Lambda}(a-a^{-1})h_{D_L}h_{D_R}.
 \label{un:h-skein}
\end{equation}
Here $\Lambda$ is the signed linking number of the smoothed components;
it need not be positive after earlier switches.  The local diagram
relations are
\begin{equation}
\begin{aligned}
 w(D^{\rm sw})&=w(D)-2,
       &r(D^{\rm sw})&=r(D),\\
 w(D)-1&=w(D_L)+w(D_R)+2\Lambda,
       &r(D)&=r(D_L)+r(D_R).
\end{aligned}
 \label{un:frame-relations}
\end{equation}
The second writhe equation counts every crossing except the smoothed
one; the last equation is additivity of tangent rotation under oriented
smoothing.  Solve \eqref{un:h-skein} for $h_D$ and multiply by
$a^{w(D)+r(D)}(a-a^{-1})$.  The switched term becomes
$\mathcal B_{D^{\rm sw}}$ by the first line of
\eqref{un:frame-relations}.  In the product term the remaining power of
$a$ is $w(D)+r(D)-1-2\Lambda$, which the second line identifies with
$w(D_L)+r(D_L)+w(D_R)+r(D_R)$.  Thus
\begin{equation}
 \mathcal B_D
   =\mathcal B_{D^{\rm sw}}+
       \mathcal B_{D_L}\mathcal B_{D_R}.
 \label{un:B-skein}
\end{equation}
This is an exact Laurent-polynomial identity with all linking and framing
factors included.

\subsection{The fixed exposed basepoint and the processed prefix}
\label{un:prefix}

Start the descending knot traversal at a point on the exposed supporting
edge, with the occupied side to its left.  At each first-under crossing
use \eqref{un:B-skein}: the switched branch proceeds along the same
component, and the smoothing branch detaches the interval from its first
visit to its second.  Rebase the detached child at an exposed point of
its own.  Keep the parent's \emph{original} exposed basepoint.

Three properties make this a closed recursive class.  First, a
first-under smoothing joins the child's end to its beginning with a
left corner.  All its other corners lie on the unprocessed original
arc and turn left.  Second, an earlier switched crossing has its first
visit in the processed prefix.  It cannot have both visits in a later
detached child, so every retained child self-crossing is still positive.
Third, the right corner created on the parent at an earlier smoothing
lies in that parent's processed prefix.  It likewise cannot enter a
future detached child.  Thus every detached child is again a uniform
positive polygon, with fewer retained self-crossings.

The actual parent diagram still contains its processed prefix, its
switched crossings, and its artificial parent corners.  The open
GLSSW current forgets that prefix only as a device for traversing the
remaining arc.  In particular, we do \emph{not} assert that the knot
obtained by closing the active arc with a straight segment is the actual
parent knot.  The objects being compared are the two recursive
calculations, with their original parent framings retained on the knot
side.

At termination every self-crossing of the actual parent is first-over,
so it is a descending unknot.  Its framed value is nevertheless to be
computed.  Let $I_b$ be the sum of the winding numbers of the two regions
adjacent to its fixed basepoint.  Whitney's formula in our determinant
convention is
\begin{equation}
 \sum_c\operatorname{sgn}\det(t_{\rm first}(c),t_{\rm second}(c))
       =I_b-r(D).
 \label{un:Whitney}
\end{equation}
For a descending diagram the first strand is the over-strand, so the
left side is its writhe.  The original exposed basepoint has adjacent
indices $0$ and $1$.  Removing a detached child does not change either
index there: that child avoids the supporting edge and has winding zero
in a neighborhood of the basepoint.  Thus $I_b=1$ for every terminal
parent.  The needed exposed-basepoint form of Whitney's identity also
has a direct local verification.  Deform the return arc within the
supporting half-plane to a simple arc, fixing its basepoint germ.
In a generic deformation, a small curl changes the rotation by
$\varepsilon\in\{1,-1\}$ and its first-visit determinant sum by
$-\varepsilon$.  A self-tangency creates two crossings with opposite
first-visit signs and changes no net rotation.  A triple-point passage
changes neither quantity.  Thus their sum stays constant; on the
positively oriented embedded circle it is one.  This proves precisely
the instance of \eqref{un:Whitney} used here and gives
\begin{equation}
 w(D_{\rm terminal})+r(D_{\rm terminal})=1.
 \label{un:terminal-frame}
\end{equation}
Since its HOMFLY polynomial is $1$, it follows that
\begin{equation}
 \mathcal B_{D_{\rm terminal}}=a(a-a^{-1})=y.
 \label{un:B-terminal}
\end{equation}

The recursion is finite.  Along any continuing parent branch there are
only finitely many unprocessed first-under decisions.  A smoothing also
lowers the retained crossing number of each detached child.  Starting
with the terminal value \eqref{un:B-terminal}, repeated use of
\eqref{un:B-skein} proves that every resulting $\mathcal B_D$ is a
polynomial in $y$, divisible by $y$.  In particular, for an initial
uniformly left-turning $Q$, there is a polynomial $F_Q(y)$ such that
\begin{equation}
 \mathcal B_Q=yF_Q(y),\qquad F_Q(y)\in\mathbb Z[y].
 \label{un:F-polynomial}
\end{equation}

For this initial diagram $w(Q)=m_Q$ and $r(Q)=R_Q$.  Substituting into
\eqref{un:framed-definition} and canceling $y=a(a-a^{-1})$ yields
\begin{equation}
 a^{m_Q+R_Q-1}h_Q(a)=F_Q(a^2-1).
 \label{un:row-polynomial}
\end{equation}
Consequently the $z^0$ row has no $a$ power below
$d_Q=1-m_Q-R_Q$, and its coefficient there is
\begin{equation}
 c(Q)=F_Q(-1)=-\mathcal B_Q\big|_{a=0}.
 \label{un:corner-evaluation}
\end{equation}
The specialization is legitimate because \eqref{un:F-polynomial} was
proved first.  No assertion about the lowest $a$ degree of higher
$z$ rows is needed.

\subsection{Identification by induction}
\label{un:identification}

\begin{proposition}
\label{un:uniform-proposition}
For every uniformly turning polygon $Q$ in the generic domain,
\begin{equation}
 \Gamma_A(Q)=(-1)^{\ell(Q)}c(Q).
 \label{un:uniform-result}
\end{equation}
\end{proposition}

\begin{proof}
For a uniformly left-turning polygon, initialize the current on an
exposed supporting edge as in \eqref{un:initial}.  Its initial normalized
value is $\nu_A(Q)$.

Compare the current traversal with the framed knot traversal just
described.  First-over encounters leave both calculations unchanged.
Second visits of switched crossings, and crossings with a removed child,
are events on neither side: the knot traversal acts only at first visits,
and such crossings are not self-crossings of the active arc.
Original left-corner passages leave $f$ unchanged and have no knot
skein operation.  At a first-under crossing, put
$g_D=-\mathcal B_D|_{a=0}$.  By the already established polynomiality
this is well-defined on every recursive branch.  Equation
\eqref{un:B-skein} gives
\begin{equation}
 g_D=g_{D^{\rm sw}}-g_{D_L}g_{D_R},\qquad
 g_{\rm terminal}=1.
 \label{un:g-recursion}
\end{equation}
For an initial uniform child, \eqref{un:corner-evaluation} identifies
$g_{D_L}=c(Q_c)$.  The child has fewer retained self-crossings, so
induction identifies this number with $\nu_A(Q_c)$.  Equations
\eqref{un:g-recursion} and \eqref{un:f-recursion} then have the same
coefficient, the same continuing parent branch, and the same terminal
value.  Induction backwards along the finite cursor recursion gives
$f=g$ at its initial state.

For the zero-crossing base case, there are no splitting decisions:
both traversals end at their terminal value $1$, and the polygon is an
unknot of rotation one.  This starts the induction on retained
self-crossings without using an amplitude anchor or a wall law.
For general $Q$, the initial equality is
$\nu_A(Q)=g_Q=c(Q)$, proving \eqref{un:uniform-result} in the
left-turning orientation.

Finally, reverse a uniformly right-turning polygon. Reverse the order at
every tree vertex. Every near and far sign changes sign, so each linear
gate changes sign. For a tree with $n-1$ leaves the sum of
$(\text{arity}-1)$ over all internal vertices is $n-2$; this is exactly
its number of gates. The number of ordinary-vertex minus signs is
unchanged. Thus the finite tree sum itself gives
$A(\overline P)=(-1)^{n-2}A(P)=(-1)^nA(P)$.
A support with $s$ cuts has
$\sum_Qv(Q)=v(P)+2s$, so its product acquires exactly $(-1)^{v(P)}$.
Therefore $\Gamma_A(\overline P)=(-1)^{v(P)}\Gamma_A(P)$.
Since $\ell(\overline P)=v(P)-\ell(P)$, the normalized $\nu_A$ is unchanged.
Reversing all strands preserves the positive lift's HOMFLY polynomial,
its retained crossing number, and $|r|$, and hence preserves $c$.
The left-turning result proves the right-turning case.
\end{proof}

% ---- completion ----
\section{Reconstruction of the amplitude}
\label{sec:completion}

The uniform calculation establishes
\begin{equation}
 \Gamma_A(Q)=(-1)^{\ell(Q)}c(Q)
 \qquad(Q\ \text{uniform}).
 \label{eq:uniform-identification-final}
\end{equation}
Together with Proposition~\ref{prop:nonuniform}, this identifies every
primitive needed in Lemma~\ref{lem:inversion}.

\begin{proof}[Proof of Theorem~\ref{thm:main}]
Apply \eqref{eq:inversion} with \(X=A\). Any decomposition containing a
nonuniform component has a zero primitive factor, so only
\(\Dcal_{\mathrm u}(P)\) remains. Substituting
\eqref{eq:uniform-identification-final} gives
\begin{equation}
 A(P)=\sum_{S\in\Dcal_{\mathrm u}(P)}
       (-1)^{\sum_{Q\in P_S}\ell(Q)}\prod_{Q\in P_S}c(Q).
 \label{eq:reconstruction-before-sign}
\end{equation}
Equation~\eqref{eq:left-count} then gives
\begin{equation}
 A(P)=(-1)^{\ell(P)}
       \sum_{S\in\Dcal_{\mathrm u}(P)}(-1)^{|S|}
                         \prod_{Q\in P_S}c(Q).
 \label{eq:reconstruction-after-sign}
\end{equation}
The right side is the definition \eqref{eq:state-sum} of \(C(P)\).
\end{proof}

\paragraph{Continuation to the whole generic locus.}
Theorem~\ref{thm:main} has been proved on the auxiliary generic domain of
Section~\ref{sec:statement}. A \emph{chamber} is a connected component of
the generic locus. Within a chamber the auxiliary domain is the complement
of finitely many proper algebraic hypersurfaces, hence open and dense: an
auxiliary determinant vanishing identically on the open chamber would
vanish identically as a rational function of the vertices, and the only
identically vanishing determinants among vertices and crossing points are
those of three points of one edge line, which never occur, since a
component of a separating smoothing contains at most one segment of each
original edge (Section~\ref{un:open}). Both sides of the identity are
constant on chambers. The tree coefficient \(A(P)\) is a function of the
signs \(H(i,k,j)\) alone, by \eqref{ep:tree-recursion}, and these are
constant along any path in a chamber, since (G1) holds throughout. The
state sum \(C(P)\) is constant on a chamber because along such a path no
vertex meets a nonincident edge and no turn passes through \(0\) or
\(\pi\): no crossing is created or destroyed, the Gauss chords and hence
the separating supports are fixed, each positive diagram moves by a
planar isotopy, and the integers \(m_Q\), \(r_Q\), \(\ell(P)\) are
unchanged. An identity between two locally constant functions that holds
on a dense subset of a chamber holds on the chamber. Hence \(A(P)=C(P)\)
on the whole generic locus. Flat turns, vertex--edge collisions, and
triple crossings lie outside the generic locus and are not treated here.
The zero-source remark after Theorem~\ref{ep:mixed-edge} is not needed
for this extension.

\paragraph{What the vertices are.}
The uniform knot weight is an effective marked vertex obtained after
resumming the finite GLSSW trees and subtracting their separating sewings.
It is not identified with an individual bare \(V^+\) or \(V^-\).
The mark makes intermediate open currents well defined. At the end,
GLSSW cyclic invariance removes the dependence on which external edge was
chosen as the marked, single-minus root. Separating smoothings reconstruct
the original closed amplitude by the finite inverse
\eqref{eq:inversion}; this is the precise algebraic content of the
disconnected-to-connected description.

\paragraph{Scope of the coefficient statement.}
The argument identifies the prescribed scalar corner. The polynomial
descending expansion used in its proof is a calculation of the framed
\(z^0\) row of each uniform knot. It does not establish a canonical
two-variable deformation of the full GLSSW amplitude, nor identify a
single Legendrian ruling polynomial with every corner. No such additional
claim is required for \eqref{eq:reconstruction-after-sign}.

% ---- physics ----
\section{Wilson loops and a possible string field theory}
\label{sec:physics}

The algebraic sewing argument has a precise Wilson-loop reformulation.
Its interpretation as the localization of a physical string field theory
is a further, \textbf{CONJECTURAL}, statement. We separate these claims.
The positive lift, plane orientation and contact-form convention are those
already fixed in the polygon construction; no additional mirror operation
is introduced here. For translation to the letter, the corner and state
sum are \texttt{[eq::factor]} and \texttt{[eq::C]}, and the HOMFLY
normalizations are \texttt{[eq::skein]} and \texttt{[eq::skein2]}.

\subsection{A normalized, framed Wilson observable}

Let $S$ be a uniform separating decomposition, with $e=|S|$ and $k=e+1$
components. Write $L_S$ for the actual smoothed link, including its mixed
crossings, and set
\begin{equation}
 p_i=m_{Q_i}+|r_{Q_i}|=1-d_{Q_i},\qquad
 p_S=\sum_{i=1}^k p_i,\qquad
 \Lambda_S=\sum_{i<j}\operatorname{lk}(Q_i,Q_j).
 \label{phys:framing-data}
\end{equation}
The $m_{Q_i}$ count self-crossings only. The mixed crossings contribute
to $\Lambda_S$, not to $p_i$.

For a Hermitian $U(N)$ connection $\mathscr A$, define
\begin{equation}
 U(Q)=\mathcal P\exp\!\left(i\oint_Q\mathscr A\right),\qquad
 W(Q)=\operatorname{Tr}_{\square}U(Q),\qquad
 W(L_S)=\prod_{i=1}^k W(Q_i).
 \label{phys:Wilson-operator}
\end{equation}
The expectation is divided by the vacuum partition function. The
fundamental trace is not divided by $N$ or by the unknot expectation.
Use the analytic Chern--Simons variables
\begin{equation}
 g_s=\frac{2\pi i}{\kappa+N},\qquad q=e^{g_s},\qquad
 a=q^{N/2},\qquad z=q^{1/2}-q^{-1/2}=2\sinh(g_s/2).
 \label{phys:variables}
\end{equation}
Here $\kappa$ is the level, while $k$ counts link components. The
large-$N$ expansion keeps $a$ fixed as $g_s\to0$.

Zero framing means that each component is framed by its Seifert
longitude. In the convention of Labastida--Mari\~no--Vafa,
Eq.~(4.54), the normalized fundamental expectation is
\begin{equation}
 \langle W(L_S)\rangle_0
 =a^{2\Lambda_S}\frac{a-a^{-1}}{z}H_{L_S}(a,z).
 \label{phys:Wilson-zero}
\end{equation}
The factor $a^{2\Lambda_S}$ is part of this link normalization
\cite{LMV}. There is only one displayed unknot factor: for a split
$k$-component link the remaining $k-1$ such factors already occur in
$H_{L_S}$.

Choose the polygon ribbon to have framing $p_i$ on $Q_i$. Concretely,
one can start with the blackboard ribbon of the positive diagram, whose
framing is $m_{Q_i}$, and add $|r_{Q_i}|$ positive full twists. We fix
the fundamental $U(N)$ framing convention to have phase
$\exp(p_i N g_s/2)=a^{p_i}$: the quadratic Casimir of the fundamental
is $N$ in this convention. Hence
\begin{equation}
 \langle W(L_S)\rangle_{\mathrm{rib}}
 =a^{p_S}\langle W(L_S)\rangle_0
 =a^{p_S+2\Lambda_S}\frac{a-a^{-1}}{z}H_{L_S}(a,z).
 \label{phys:Wilson-ribbon}
\end{equation}
This specifies the ribbon completely as a framed-knot class. It does
not identify its framing integer with the signed transverse
self-linking number: the degree-bound convention gives
$\operatorname{sl}(T_Q)=-m_Q-|r_Q|=-p_Q$. These are different signed
quantities. Nor is a ribbon the conormal brane introduced below.
Changing the ribbon changes the regularization of the quantum Wilson
expectation, although its central contour and classical holonomy remain
the same. The mutual-linking correction in
\eqref{phys:Wilson-zero} must not be added to $p_S$ a second time.

The lowest HOMFLY coefficient of a $k$-component link factorizes as
\begin{equation}
 [z^{1-k}]H_{L_S}
 =a^{-2\Lambda_S}(a-a^{-1})^{k-1}
   \prod_i [z^0]H_{Q_i}.
 \label{phys:lowest-H}
\end{equation}
This is the Lickorish--Millett relation, in the form of
\cite[Eq.~(4.64)]{LMV}. Inserting it into
\eqref{phys:Wilson-zero} cancels $a^{2\Lambda_S}$ and supplies one
additional factor of $(a-a^{-1})/z$. Thus
\begin{equation}
 [z^{-k}]\langle W(L_S)\rangle_0
 =\prod_i D_{Q_i}(a),\qquad
 D_Q(a)=(a-a^{-1})[z^0]H_Q(a,z).
 \label{phys:disk-factorization}
\end{equation}
This coefficient identity removes the mixed-linking dependence at the
selected order; it does not unlink the original link.

The corner floor $\min\deg_a[z^0]H_Q\ge d_Q=1-p_Q$ implies
$\min\deg_a D_Q\ge-p_Q$. At degree $-p_Q$, the $+a$ term in
$a-a^{-1}$ would require degree $d_Q-2$ in $[z^0]H_Q$, below the floor.
The $-a^{-1}$ term gives $-c(Q)$. Consequently
\begin{equation}
 [a^{-p_Q}]D_Q=-c(Q),\qquad
 [a^0z^{-k}]\langle W(L_S)\rangle_{\mathrm{rib}}
 =(-1)^k\prod_i c(Q_i).
 \label{phys:Wilson-corner}
\end{equation}
In the product, the sum of the floors can only be attained when every
factor is at its own floor; if any corner vanishes, so does this
coefficient.

An exact full-link extension of the corner sum is therefore
\begin{equation}
 \boxed{\displaystyle
 \mathcal C_P^{W}(a,z)=(-1)^{\ell(P)+1}
   \sum_{S\ \mathrm{uniform}}z^{k_S}
       \langle W(L_S)\rangle_{\mathrm{rib}},\qquad
 [a^0]\bigl([z^0]\mathcal C_P^{W}\bigr)=C(P).}
 \label{phys:observable}
\end{equation}
Indeed $k_S=|S|+1$, so the sign supplied by
\eqref{phys:Wilson-corner} is
$(-1)^{\ell(P)+1+k_S}=(-1)^{\ell(P)+|S|}$, precisely the state-sum
sign. There is no additional $(-1)^{|S|}$ in
\eqref{phys:observable}; including it would count that sign twice.
This is a finite linear combination of genuine framed Wilson
correlators. Its nonzero-$a,z$ dependence is a chosen extension of the
corner observable. The equality proved here concerns its indicated
coefficient, not a universal finite-parameter deformation of GLSSW.
For a general full link, component corner bounds alone also do not
justify replacing the coefficient extraction by a simultaneous
numerical limit at $(a,z)=(0,0)$.

\subsection{What the disk sector means}

The A-model is the topological string obtained from a supersymmetric
sigma model by a twist that localizes its amplitudes on
pseudoholomorphic maps into a symplectic target. Open worldsheets have
boundary on Lagrangian branes. With $N$ branes on the zero section of
$T^*S^3$, its open-string field theory gives Chern--Simons theory on
$S^3$. The large-$N$ conifold transition replaces
the compact branes by the resolved geometry
$X=\mathcal O(-1)\oplus\mathcal O(-1)\to\mathbb P^1$; the dictionary
used here is reviewed in \cite[Sec.~2]{LMV}.

A knot $Q\subset S^3$ supplies the conormal Lagrangian
\begin{equation}
 L_Q=N^*Q
 =\{(x,\xi)\in T^*S^3:x\in Q,\ \xi(\dot Q)=0\}.
 \label{phys:conormal}
\end{equation}
It is three-dimensional: one coordinate along $Q$ and a
two-dimensional cotangent fiber. The appropriate displaced or continued
brane survives on the resolved side. The worldsheet boundary lies on
this brane and has a specified winding around its distinguished circle.
It is not fixed pointwise to the original knot in the zero section.
An open disk map in the six-dimensional target therefore says nothing
about whether $Q$ bounds an embedded disk in three dimensions.

In the knot/string dictionary, strings between the compact branes and
the knot branes generate Wilson insertions. With auxiliary brane
holonomies $V_i$ one may use the source operator
\begin{equation}
 Z(U,V)=\exp\!\left(
  \sum_i\sum_{n\ge1}\frac{1}{n}
  \operatorname{Tr}\bigl(U(Q_i)^n\bigr)
  \operatorname{Tr}\bigl(V_i^n\bigr)\right).
 \label{phys:OV-source}
\end{equation}
Choosing the coefficient linear in each $\operatorname{Tr}V_i$
selects the product of fundamental Wilson insertions. Reversing the
probe-brane orientation replaces $V_i$ by its inverse; the orientation
in \eqref{phys:OV-source} is our convention. This is the
Ooguri--Vafa source operator, as reviewed in \cite[Eq.~(2.6)]{LMV}.
The logarithm of the expectation generates connected worldsheets;
the expectation itself also contains disconnected worldsheets.

A connected genus-$g$ surface with $h$ boundaries carries
$g_s^{2g-2+h}$. In the selected source sector there are $k$ boundaries,
each with winding one. For $b$ connected open components of total genus
$G$, the coupling power is therefore
\begin{equation}
 g_s^{2G-2b+k},\qquad b\le k,\qquad G\ge0.
 \label{phys:worldsheet-power}
\end{equation}
The smallest exponent is $-k$, attained only for $b=k$, $G=0$: a
disconnected collection of $k$ genus-zero one-boundary worldsheets,
one per source. This is the meaning of ``one disk per knot.'' The disk
block is a sum of maps and their weights; the knot itself is never
declared to be a disk. Since $z=g_s+O(g_s^3)$, the leading $z^{-k}$
coefficient has the same interpretation. Equation
\eqref{phys:disk-factorization} is its exact knot-polynomial
realization. Connecting two of the fixed Wilson boundaries into an
annulus raises the coupling exponent by two. Such interactions may
retain linking information, and occur beyond the selected order.

\subsection{Framing, boundary holonomy and the selected charge}

In zero framing the coefficient used by the state sum is
$[a^{-p_S}z^{-k}]\langle W(L_S)\rangle_0$. If a term is written
$(a^2)^{\beta}$, it selects total exponent
\begin{equation}
 \beta=-\frac{p_S}{2}.
 \label{phys:selected-charge}
\end{equation}
The ribbon shifts this to $\beta+p_S/2=0$. In an open-string expansion
the exponent labels a relative charge with a chosen boundary
trivialization. It need not be a nonnegative integer, and is not by
itself the area of an embedded surface.

The framing shift can equivalently be placed in the boundary source.
For the fundamental, single-winding term,
\begin{equation}
 \bigl(a^{p_i}\langle W(Q_i)\rangle_0\bigr)
        \operatorname{Tr}V_i
 =\langle W(Q_i)\rangle_0
        \operatorname{Tr}(a^{p_i}V_i).
 \label{phys:source-shift}
\end{equation}
Thus an extra power of $a$ assigned to the holonomy may simply be the
same framing factor expressed in a different boundary coordinate. It
must not be included twice. If a flat boundary coordinate is further
redefined by a framing-dependent sign, that sign belongs to the source
convention; it is separate from our stipulated $a^{p_i}$ Wilson framing
phase.

With $t=Ng_s$, our variables give $a=e^{t/2}$. The conventional
large-volume closed coordinate is $e^{-t}=a^{-2}$. Accordingly the
formal lowest-$a$ coefficient does not, merely by notation, establish a
large-volume limit at $a\to0$. Such a physical limit would require an
analytic chamber and boundary-coordinate prescription. The coefficient
identity \eqref{phys:observable} needs neither that extra assertion nor
a change of the fixed contact convention.

\subsection{The disconnected prescription and its limitations}

The factor in \eqref{phys:observable} is $z^{k}=z\,z^{e}$. One factor
of $z$ removes a common disk pole; the remaining $z^e$ has a natural
tree-sewing power count. Start with $k$ abstract disks, of total Euler
characteristic $k$. Attaching a boundary strip lowers Euler
characteristic by one. After $e$ strips,
\begin{equation}
 \chi=k-e,\qquad g_s^{-\chi}=g_s^{e-k}.
 \label{phys:strips}
\end{equation}
For a connected tree $e=k-1$, so the resulting surface has disk order
$g_s^{-1}$. Reverse the separating smoothings of a decomposition: each
step joins two previously distinct components, hence the graph of
joins is a tree. Conversely, a join within a previously connected
collection creates a graph cycle. Depending on the ribbon attachment,
such a cycle can create another boundary or a handle; it should not
automatically be called positive genus.

This prospective sewing changes the boundary decomposition. It is
different from the annulus interaction with the original $k$ Wilson
sources held fixed. Multiplying a correlator by $z^e$ alone performs no
geometric sewing and supplies no integration over attachment data.
Euler characteristic explains the leading $g_s^e$ cost of strips, but
does not determine the exact function $z^e$, its sign or its framing
weight. An exact string construction would need to reproduce the skein
factor, rather than infer it solely from the leading genus power.

There is a useful precedent in the disconnected twistor prescription.
The CSW rules join MHV vertices by propagators into trees.
Gukov--Motl--Neitzke compare descriptions by disconnected
holomorphic curves with a connected description: propagator poles can,
under their contour assumptions, localize attachment points to nodes
\cite[Secs.~2.2 and 4]{GMN}. Their curves belong to the twistor B-model,
whereas the knot disk blocks above belong to an open A-model dictionary.
The nodal combinatorics supplies a precedent, not an identification of
the targets, degrees or propagators. Their equivalence of prescriptions
also does not by itself fix our sum over all component counts.

The present algebra gives more specific data for such a comparison.
The uniform block is the smoothing-irreducible quantity
$\Gamma_A(Q)=(-1)^{\ell(Q)}c(Q)$. It is obtained by regrouping GLSSW
trees. It is not a single bare $V^+$ or $V^-$ interaction. A marked
attachment survives in the open current; closing that current too early
would discard the data needed for the sewing identity. Thus a proposed
``single-minus disk vertex'' should mean this effective marked block,
with the external single-minus label selecting the overall root.
Assigning an independent physical negative-helicity particle to every
component is not established by the regrouping.

\subsection{What would make this a physical construction?}

A physical realization would have to provide the following data.
\begin{enumerate}
\item A state space on marked boundary intervals, including the
crossing and framing data retained by the algebraic current. Its
relation to the half-collinear GLSSW external states must be explicit.
\item A kinetic operator and a gauge-fixed inverse that propagates
those states along a strip. In an open-string description the
propagating objects are open-string states between the relevant branes;
in a Chern--Simons reduction a gauge connection represents the
appropriate string-field sector. Neither description identifies the
strip with an ordinary four-dimensional gluon propagator without a
further reduction.
\item A pairing and gluing measure, with integration contours or a
localization rule, that turn the propagator into the precise marked
crossing kernel. This must derive the separating restriction, signs,
multiplicities and exact $z$ dependence.
\item A BRST or equivalent gauge structure that explains cyclic/root
independence and makes the sewn observable independent of auxiliary
choices. A finite matrix identity alone is not a kinetic operator or
a BRST construction.
\item Boundary holonomies and a charge projection reproducing
\eqref{phys:Wilson-corner}, including the ribbon integers $p_i$, as well
as a reason the physical localization has the uniform effective blocks
of the algebraic argument.
\end{enumerate}
These requirements specify a \textbf{CONJECTURAL} open-string or
crossing-defect theory. The direct combinatorial sewing supplies its
candidate effective vertices and exact target observable; the standard
Chern--Simons/A-model dictionary supplies the Wilson and disk
interpretations. Neither, on its own, constructs the required strip
propagator.

% ---- references ----
\begin{thebibliography}{3}
\bibitem{GLSSW}
A.~Guevara, A.~Lupsasca, D.~Skinner, A.~Strominger and K.~Weil,
\emph{Single-minus gluon tree amplitudes are nonzero},
\href{https://arxiv.org/abs/2602.12176}{arXiv:2602.12176} (2026).

\bibitem{LMV}
J.~M.~F.~Labastida, M.~Mari\~no and C.~Vafa,
\emph{Knots, links and branes at large $N$},
JHEP \textbf{11} (2000) 007,
\href{https://arxiv.org/abs/hep-th/0010102}{arXiv:hep-th/0010102}.
The fundamental Wilson normalization and lowest-link relation are
Eqs.~(4.54) and (4.64); the Wilson-source construction is Eq.~(2.6).

\bibitem{GMN}
S.~Gukov, L.~Motl and A.~Neitzke,
\emph{Equivalence of twistor prescriptions for super Yang--Mills},
\href{https://arxiv.org/abs/hep-th/0404085}{arXiv:hep-th/0404085}.
See Sections~2.2 and~4 for disconnected curves, propagators and
localization to nodes, including the contour and normalization qualifications.
\end{thebibliography}
\end{document}
